% Generated draft: requires mathematical and source review.
\documentclass[11pt]{article}
\usepackage[margin=25mm]{geometry}
\usepackage{fontspec}
\setmainfont{Latin Modern Roman}
\usepackage{amsmath,amssymb,mathtools,unicode-math}
\setmathfont{Latin Modern Math}
\usepackage{enumitem,needspace,xurl,hyperref}
\setlength{\parindent}{0pt}
\setlength{\parskip}{5pt}
\setlength{\emergencystretch}{3em}
\newcommand{\N}{\mathbb{N}}
\newcommand{\Z}{\mathbb{Z}}
\newcommand{\Q}{\mathbb{Q}}
\newcommand{\R}{\mathbb{R}}
\newcommand{\C}{\mathbb{C}}
\newcommand{\F}{\mathbb{F}}
\newcommand{\A}{\mathbb{A}}
\newcommand{\PP}{\mathbb{P}}
\newcommand{\E}{\mathbb{E}}
\newcommand{\eps}{\varepsilon}
\newcommand{\dd}{\,\mathrm{d}}
\begin{document}
\section*{UnsolvedMath ID: 20002330}
\textbf{Source problem code:} catalogue entry
\subsection{Definitions and mathematical statement}
Problem 15 asks to study the group configuration theorem in continuous logic. Its trailing ``Background:'' has no supplied continuation. The problem does not fix stability/simplicity hypotheses, a configuration definition, parameters, or the desired notion of interpretable group.

The aid reports an existing transfer from simple theories to simple thick compact abstract theories (CATs): in the CAT presentation of a simple continuous theory, an algebraic quadrangle yields a nontrivial gradedly almost hyperdefinable group, with the associated action in the relevant construction. The full quadrangle independence/algebraicity requirements and meanings of the graded hyperdefinability conventions are not reproduced, so no unconditional group theorem for every continuous theory is asserted.

Its candidate reduction addresses a further, explicitly conditional metric upgrade. Suppose the CAT-level group and action have presentations
\[
 G_0=D_G/\{\rho_G=0\},\qquad X_0=D_X/\{\rho_X=0\},
\]
where the pseudometrics and presentation domains are definable in the sense required to form ordinary continuous imaginary sorts. Require the zero-distance relations to represent the intended equality relations, not an arbitrary further collapse. Suppose multiplication, inverse, identity, and the action are well-defined definable maps in these presentations and have supplied uniform continuity moduli; in particular multiplication and action must be uniformly continuous jointly on their product domains, not merely separately continuous.

Let \(\widehat G\) and \(\widehat X\) be the metric completions. The aid reports unique continuous extensions
\[
 \widehat m:\widehat G\times\widehat G\to\widehat G,
 \quad\widehat{\mathrm{inv}}:\widehat G\to\widehat G,
 \quad\widehat a:\widehat G\times\widehat X\to\widehat X,
\]
with the group and action identities extending from the dense quotients. With the assumed definable-pseudometric sort presentation, these give an interpretable complete metric group and action. Uniqueness refers to the extensions of the fixed presented maps. Metric completion by itself does not turn arbitrary hyperdefinable domains or maps into ordinary definable ones.

Thus the remaining upgrade task in this route is to construct faithful definable metric presentations and uniform moduli for the CAT output. The aid does not prove such presentations always exist, nor identify an unrestricted ordinary-continuous-imaginary group configuration theorem. Its result concerns sufficiency once those extra data are available.
\subsection{Short English statement}
The source proposes a continuous-logic group-configuration program. The aid reports a CAT-level group construction and shows that a suitable definable, uniformly continuous metric presentation extends through quotient and completion. It leaves construction of that presentation as the actual additional requirement for an ordinary metric interpretation.
\subsection{Sources}
\begin{itemize}
\item Local UnsolvedMath registry, permanent ID 20002330, exact source code catalogue entry; canonical \texttt{aim-logic-notes.json}, \texttt{notes[105]}, tag \texttt{problem}; AIM Workshop Problem Lists, Logic, \emph{Model theory of metric structures}, Problem 15. Created and updated 14 August 2026.
\item Source URL \url{https://aimath.org/WWN/continuouslogic/continuouslogic.pdf}. The extracted background is empty after its heading. The PDF was not accessed.
\item The aid refers to the group-configuration transfer for simple thick CATs without naming an author, publication, theorem number, or precise quadrangle statement. No authoritative full citation is supplied for that transfer or the metric criterion.
\item Attempt 1: \texttt{partial\_result}, problem status at run \texttt{partially\_solved}, rigor label \texttt{proved}, novelty confidence \texttt{low}; classification \texttt{PARTIAL-PROGRESS}, registry status \texttt{partially\_solved}. Candidate: definable-pseudometric presentation and completion criterion. These are dated supplied assessments, not an unrestricted solution.
\item Ulam AI / UnsolvedMath Contributors; CC BY 4.0. Machine-generated research aid, not peer reviewed; independently verify before citation.
\end{itemize}
\end{document}
